\expandafter\def\csname HistoricalNoteText15\endcsname{Dass sie immer einen ganz bestimmten Verlauf hat, beruht darauf, dass sie einen Teil der äußeren Begrenzung einer Fläche bilden soll, für die $T$ ein bestimmtes Vorzeichen hat. Ich habe das positive Vorzeichen gewählt, was an sich völlig beliebig ist. So verstanden setzt sich beispielsweise die bei $1$ eintretende Linie durch $11$ nach $10$ fort. Als Teil der Grenzlinie einer Fläche, in der $T$ negativ ist, müsste die Linie $1.11$ nach $2$ fortgesetzt werden. Spricht man dagegen nur von einer Linie, in der $T=0$ ist, ohne sie als Teil der Begrenzung einer bestimmten Fläche zu betrachten, so könnte eher $11.9$ als natürliche Fortsetzung von $1.11$ gelten. Der hier gewählte Gesichtspunkt unterscheidet mein gegenwärtiges Verfahren von dem von $1799$ und trägt wesentlich zur Vereinfachung der Beweisführung bei.}
